\section{Conclusion}

This paper introduced the Physics-Informed Market Strategy (PIMS) framework, bridging micro-level behavioral memory strategies with macro-level non-equilibrium swarm physics on the Indian Nifty 50 TRI Index (2015 to 2025).

In Stage 1, we established that standalone memory kernels require rigorous out-of-sample validation, with microscopic queue-level sentiment injection achieving superior risk-adjusted performance (Sharpe 1.486 for Opportunist injected with Traditionalist, and Sharpe 1.470 with 15.2\% turnover for Pushover injected with Contrarian).

In Stage 2, we derived the mean-field kinetic equations of motion in the thermodynamic limit, proving analytically and empirically that peer sentiment interaction ($p_{\text{peer}}=1.0$) induces a supercritical pitchfork bifurcation at critical memory depth $M_c = \frac{\pi}{2(1-2\eta)^2} \approx 9.82$ trading days. Isolated price observation ($p_{\text{peer}}=0$) remains strictly ergodic ($k > 0$), establishing that social peer contagion is a mandatory condition for collective financial herding. Furthermore, panic relaxation experiments demonstrated critical slowing down in conformists, with structural macro trend anchors ($\phi_C \ge 10\%$) acting as circuit breakers against permanent sentiment traps.

Finally, we proved that quantitative backtest overfitting is the empirical manifestation of the crowd pitchfork phase transition. Applying the theoretical PIMS physical stability constraint ($M_{\text{eff}} \le M_c$) as a pre-fit prior filter eliminates overfitted configurations prior to backtesting, boosting mean out-of-sample Sharpe ratios by +41.4\% while reducing the overfitting gap by 37.7\%.
